\section{Monoidal analog of cluster-tilting theory}
\label{Sec : monoidal analog of cluster tilting theory}

Let $\mathscr{C}$ and $N$ be as in
Subsection~\ref{ssec : settings in this paper}.

\subsection{Maximal real sets}
\label{ssec : maximal real set}

\begin{Def}
\label{Def : maximal real set}
Let $X$ be a set.
A set $\{L_\lambda\}_{\lambda \in X} \subset \Irr\mathscr{C}$
is called a \emph{maximal real set}
if it satisfies the following conditions:
\begin{enumerate}
    \item[(i)] $L_\lambda$ is a real prime module
    for all $\lambda \in X$,
    \item[(ii)] $\mathfrak{d}(L_{\lambda_1}, L_{\lambda_2}) = 0$
    for all $\lambda_1, \lambda_2 \in X$,
    \item[(iii)] if a real prime module $L \in \Irr\mathscr{C}$
    satisfies $\mathfrak{d}(L_\lambda, L) = 0$
    for all $\lambda \in X$,
    then $L \in \{L_\lambda\}_{\lambda \in X}$.
\end{enumerate}
\end{Def}

This is a monoidal analog of a maximal rigid object. 

Note that for any commuting family $\mathcal{S}$ of real prime modules,
there is a maximal real set which contains $\mathcal{S}$
by Zorn's lemma.

\begin{Lem}
\label{Lem : finiteness of maximal real sets}
A maximal real set in $\mathscr{C}$ is a finite set.
\end{Lem}

\begin{proof}
Let $\{L_\lambda\}_{\lambda \in X}$ be a maximal real set.
We choose $\lambda_1, \ldots, \lambda_k \in X$ such that
$\{\mathrm{PBW}(L_{\lambda_i}) \mid 1 \leq i \leq k\}$
forms a basis of
$\sum_{\lambda \in X}\Q\mathrm{PBW}(L_\lambda) \subset \Q^N$.
For $\lambda \in X$, we can write
\[
m\mathrm{PBW}(L_\lambda)
+ \sum_{i=1}^k m_i\mathrm{PBW}(L_{\lambda_i})
= \sum_{i=1}^k p_i\mathrm{PBW}(L_{\lambda_i}),
\]
where $m \in \Z_{>0}$, $m_i, p_i \in \Z_{\geq 0}$,
and $m_i p_i = 0$ for all $i$.
Since $\{L_\lambda\} \cup \{L_{\lambda_i} \mid 1 \leq i \leq k\}$
is a commuting set,
\[
L_\lambda^{\otimes m}
\otimes \bigotimes_{i=1}^k L_{\lambda_i}^{\otimes m_i},
\qquad
\bigotimes_{i=1}^k L_{\lambda_i}^{\otimes p_i}
\]
are both simple.
Then these are simple modules with the same PBW parameter,
thus we have
\begin{equation}
\label{eq : PBW relation in a maximal real set}
L_\lambda^{\otimes m}
\otimes \bigotimes_{i=1}^k L_{\lambda_i}^{\otimes m_i}
\simeq \bigotimes_{i=1}^k L_{\lambda_i}^{\otimes p_i}.
\end{equation}
Since $K(\mathscr{C})_{\Z}$ is a unique factorization domain,
$[\bigotimes_{i=1}^k L_{\lambda_i}] \in K(\mathscr{C})_{\Z}$
has a factorization into prime elements as a ring-theoretic meaning.
We write $\{q_1, \ldots, q_\ell\}$ for the set of pairwise
non-associate prime elements appearing in this factorization.
For each $\lambda \in X$, there exists a unique
$\mathbf{a}_\lambda \in \Z_{\geq 0}^{\ell}$ such that
$[L_\lambda]$ is associate to
$\prod_{j=1}^{\ell}q_j^{a_{\lambda, j}}$.
Suppose that there exist $\lambda \neq \lambda' \in X$ such that
$\mathbf{a}_\lambda - \mathbf{a}_{\lambda'}
\in \Z_{\geq 0}^{\ell}$.
Then there is some $\phi \in K(\mathscr{C})_{\Z}$ such that
$[L_\lambda] = [L_{\lambda'}]\phi$.
Since $L_{\lambda'}$ is real,
$L_\lambda \simeq L_{\lambda'} \otimes Y$
for some $Y \in \Irr\mathscr{C}$
by Proposition~\ref{Prop : divisibility by a real simple module}.
It contradicts the primeness of $L_\lambda$.
Thus $\mathbf{a}_\lambda - \mathbf{a}_{\lambda'} \notin \Z_{\geq 0}^{\ell}$ 
holds for any $\lambda \neq \lambda'$.
To finish the proof, it is enough to prove that
for any infinite sequence
$(\mathbf{a}_n)_{n \in \Z_{\geq 0}} \subset \Z_{\geq 0}^m$,
there exist $n_1 \neq n_2 \in \Z_{\geq 0}$ such that
$\mathbf{a}_{n_2} - \mathbf{a}_{n_1} \in \Z_{\geq 0}^m$.
For any infinite sequence in $\Z_{\geq 0}$,
it is easy to see that there exists a non-decreasing infinite subsequence.
In $(\mathbf{a}_n)_{n \in \Z_{\geq 0}}$,
we take an infinite subsequence whose first component is non-decreasing,
and then take an infinite subsequence whose next component
is non-decreasing.
Repeating this up to the $m$-th component, we obtain a subsequence
$(\mathbf{a}_{n_i})_{i \in \Z_{\geq 0}}$
which is non-decreasing in every component.
In particular,
$\mathbf{a}_{n_2} - \mathbf{a}_{n_1} \in \Z_{\geq 0}^m$.
\end{proof}


We can regard a maximal real set as a monoidal analog
of a cluster-tilting object by the following proposition.

\begin{Prop}
\label{Prop : simple modules commuting with a maximal real set}
Let $\{L_i \mid 1 \leq i \leq r\}$ be a maximal real set.
Let $X \in \Irr\mathscr{C}$.
If $X$ commutes with all $L_i$,
then there exists $\mathbf{a} \in \Z_{\geq 0}^r$
such that $X \simeq L^{\mathbf{a}}$.
\end{Prop}

\begin{proof}
By the definition of a maximal real set,
it is enough to show that if $X$ commutes with all $L_i$,
then $X$ is real.
Let $\{T_k \mid 1 \leq k \leq N\}$ be one of the monoidal clusters.
By Lemma~\ref{Lem : decrease of d}
and the reality of the monomials in $\{T_k\}$,
there exists $\mathbf{b} \in \Z_{\geq 0}^N$ such that
$X \nabla T^{\mathbf{b}}$ commutes with all $T_k$.
Since $\{T_k \mid 1 \leq k \leq N\}$ is a Laurent family
by Proposition~\ref{Prop : monoidal clusters are Laurent families},
there exists $\mathbf{c} \in \Z_{\geq 0}^N$ such that
$X \nabla T^{\mathbf{b}} \simeq T^{\mathbf{c}}$.
By the same argument, there exists $\mathbf{d} \in \Z_{\geq 0}^r$
such that $L^{\mathbf{d}} \nabla T^{\mathbf{b}}$ and
$L^{\mathbf{d}} \nabla T^{\mathbf{c}}$ commute with all $L_i$.
Since $L^{\mathbf{d}}$, $T^{\mathbf{b}}$, and $T^{\mathbf{c}}$
are real and
\[
\mathfrak{d}(L^{\mathbf{d}},
L^{\mathbf{d}} \nabla T^{\mathbf{b}})
= \mathfrak{d}(L^{\mathbf{d}},
L^{\mathbf{d}} \nabla T^{\mathbf{c}}) = 0,
\]
we have that $L^{\mathbf{d}} \nabla T^{\mathbf{b}}$ and
$L^{\mathbf{d}} \nabla T^{\mathbf{c}}$ are real
by Lemma~\ref{Lem : reality of head commuting with a factor}.
Then there exist $\mathbf{e}, \mathbf{f} \in \Z_{\geq 0}^r$
such that
\[
L^{\mathbf{d}} \nabla T^{\mathbf{b}} \simeq L^{\mathbf{e}},
\qquad
L^{\mathbf{d}} \nabla T^{\mathbf{c}} \simeq L^{\mathbf{f}},
\]
because $\{L_i \mid 1 \leq i \leq r\}$ is a maximal real set.
Now we consider
$L^{\mathbf{d}} \otimes X \otimes T^{\mathbf{b}} \in \mathscr{C}$.
Since $L^{\mathbf{d}} \otimes X \in \Irr\mathscr{C}$
and $T^{\mathbf{b}}$ is real,
it has a simple head by
Proposition~\ref{Prop : simplicity of heads and socles}.
There are two surjections:
\[
\begin{aligned}
L^{\mathbf{d}} \otimes X \otimes T^{\mathbf{b}}
&\twoheadrightarrow
L^{\mathbf{d}} \otimes (X \nabla T^{\mathbf{b}})
\simeq L^{\mathbf{d}} \otimes T^{\mathbf{c}}
\twoheadrightarrow L^{\mathbf{f}}, \\
L^{\mathbf{d}} \otimes X \otimes T^{\mathbf{b}}
&\simeq X \otimes L^{\mathbf{d}} \otimes T^{\mathbf{b}}
\twoheadrightarrow X \otimes L^{\mathbf{e}}.
\end{aligned}
\]
Since $L^{\mathbf{f}}$ and $X \otimes L^{\mathbf{e}}$ are simple,
both are isomorphic to the head of
$L^{\mathbf{d}} \otimes X \otimes T^{\mathbf{b}}$,
thus we have $L^{\mathbf{f}} \simeq X \otimes L^{\mathbf{e}}$.
By the reality of $L^{\mathbf{f}}$, it follows that $X$ is real.
\end{proof}

\begin{Lem}
\label{Lem : localization at a maximal real set}
Let $\{L_i \mid 1 \leq i \leq r\}$ be a maximal real set, and let
\[
G \coloneq
\langle [L_i]^{\pm 1} \mid 1 \leq i \leq r \rangle
\subset \bigl(\operatorname{Frac} K(\mathscr{C})_{\Z}\bigr)^\times.
\]
Then
\[
\Z[G] = K(\mathscr{C})_{\Z}
\bigl[[L_i]^{-1} \mid 1 \leq i \leq r\bigr]
\]
as subalgebras of $\operatorname{Frac} K(\mathscr{C})_{\Z}$.
Moreover, $G$ is a free $\Z$-module.
\end{Lem}

\begin{proof}
The inclusion
$\Z[G] \subset K(\mathscr{C})_{\Z}
\bigl[[L_i]^{-1} \mid 1 \leq i \leq r\bigr]$
is clear.
For $X \in \Irr\mathscr{C}$,
there exists $\mathbf{a} \in \Z_{\geq 0}^r$
such that any simple subquotient of $L^{\mathbf{a}} \otimes X$
commutes with all $L_i$ by Lemma~\ref{Lem : decrease of d}.
By Proposition~\ref{Prop : simple modules commuting with a maximal real set},
any simple subquotient of $L^{\mathbf{a}} \otimes X$
is a monomial in $\{L_i \mid 1 \leq i \leq r\}$.
It follows that $[X]$ is a Laurent polynomial in
$\{[L_i] \mid 1 \leq i \leq r\}$, and
$\Z[G] \supset K(\mathscr{C})_{\Z}
\bigl[[L_i]^{-1} \mid 1 \leq i \leq r\bigr]$.

Next, we prove that $G$ is free.
Since $G$ is finitely generated,
it is enough to prove that $G$ is torsion-free.
Suppose $([L^{\mathbf{a}}]/[L^{\mathbf{b}}])^k = 1$ for $\mathbf{a}, \mathbf{b} \in \Z_{\geq 0}^r$ and 
$k \in \Z_{\geq 1}$. 
Then we have $[L^{k\mathbf{a}}] = [L^{k\mathbf{b}}]$ in $K(\mathscr{C})_{\Z}$. 
Then it holds that $k\mathrm{PBW}(L^{\mathbf{a}}) = k\mathrm{PBW}(L^{\mathbf{b}})$, 
hence $L^{\mathbf{a}}$ and $L^{\mathbf{b}}$ have the same PBW parameter and they are isomorphic. 
Therefore $[L^{\mathbf{a}}]/[L^{\mathbf{b}}] = 1$ holds and 
$G$ is torsion free. 
\end{proof}
Note that $\Z[G]$ as a subalgebra of $\operatorname{Frac} K(\mathscr{C})_{\Z}$ is 
isomorphic to $\Z[G]$ as a group ring, 
that because simple classes are linearly independent in $K(\mathscr{C})_{\Z}$. 


\begin{Prop}
\label{Prop : real prime modules give prime elements}
Let $L \in \Irr\mathscr{C}$ be a real prime module.
Then $[L] \in K(\mathscr{C})_{\Z}$ is a prime element.
\end{Prop}

\begin{proof}
Since $K(\mathscr{C})_{\Z}$ is a unique factorization domain,
it is enough to prove that $[L] \in K(\mathscr{C})_{\Z}$
is irreducible.
We take $x, y \in K(\mathscr{C})_{\Z}$ such that $[L] = xy$.
Let $\{L = L_1, \allowbreak L_2, \ldots, L_r\}$
be a maximal real set,
and let $G$ be as in
Lemma~\ref{Lem : localization at a maximal real set}.
Since $[L]$ is invertible in $\Z[G]$, so are $x$ and $y$.
By Lemma~\ref{Lem : localization at a maximal real set}, we have
\[
\Z[G]^\times
= \left\{
\pm \frac{[L^{\mathbf{b}}]}{[L^{\mathbf{a}}]}
\;\middle|\; \mathbf{a}, \mathbf{b} \in \Z_{\geq 0}^r
\right\}.
\]
Therefore,
\[
x = \varepsilon_1
\frac{[L^{\mathbf{b}_1}]}{[L^{\mathbf{a}_1}]},
\qquad
y = \varepsilon_2
\frac{[L^{\mathbf{b}_2}]}{[L^{\mathbf{a}_2}]}
\]
for $\varepsilon_1, \varepsilon_2 \in \{\pm 1\}$
and $\mathbf{a}_1, \mathbf{b}_1, \mathbf{a}_2, \mathbf{b}_2
\in \Z_{\geq 0}^r$.
Then we have
$[L^{\mathbf{b}_1}] = [L^{\mathbf{a}_1}]\varepsilon_1 x$
in $K(\mathscr{C})_{\Z}$.
By Proposition~\ref{Prop : divisibility by a real simple module},
there exists $X \in \Irr\mathscr{C}$ such that
$\varepsilon_1 x = [X]$.
By the same argument, there exists $Y \in \Irr\mathscr{C}$
such that $\varepsilon_2 y = [Y]$.
Then we have $[L] = \varepsilon_1\varepsilon_2[X \otimes Y]$
in $K(\mathscr{C})_{\Z}$.
By the primeness of $L$, at least one of $X$ and $Y$ is trivial,
and we have that at least one of $x$ and $y$ is $\pm 1$.
\end{proof}

By Proposition~\ref{Prop : real prime modules give prime elements},
the following corollary is immediate.

\begin{Cor}
\label{Cor : unique prime tensor factorization in C}
For a real module $L \in \mathscr{C}$,
tensor factorization of $L$ into prime modules
is unique up to permutation.
\end{Cor}

\begin{Cor}
\label{Cor : unique prime tensor factorization in Cg}
For a real module $L \in \mathscr{C}_{\mathfrak{g}}$,
tensor factorization of $L$ into prime modules
is unique up to permutation.
\end{Cor}

\begin{proof}
By Proposition~\ref{Prop : Cg factorization},
it is enough to prove that any real module
in $\mathscr{C}_{\mathfrak{g}}^0$
has a unique tensor factorization into prime modules.
By Proposition~\ref{Prop : whole category}
and Proposition~\ref{Prop : PBW parameterization},
any module $L \in \mathscr{C}_{\mathfrak{g}}^0$
is in $\mathscr{C}_{\mathfrak{g}}^{[c, d], \mathfrak{s}}$
for some finite interval $[c, d]$
and an admissible sequence $\mathfrak{s}$.
Now the claim is clear
by Corollary~\ref{Cor : unique prime tensor factorization in C}.
\end{proof}

\begin{Lem}
\label{Lem : real prime modules give prime elements over C}
Let $L \in \Irr\mathscr{C}$ be a real prime module.
Then $[L] \in \C \otimes_{\Z} K(\mathscr{C})_{\Z}$
is a prime element.
\end{Lem}

\begin{proof}
Let $\{L = L_1, L_2, \ldots, L_r\}$
be a maximal real set in $\mathscr{C}$.
Since $\C \otimes_{\Z} K(\mathscr{C})_{\Z}$
is a polynomial ring over $\C$,
it is enough to prove that $[L]$ is an irreducible element.
We take $x, y \in \C \otimes_{\Z} K(\mathscr{C})_{\Z}$
such that $[L] = xy$.
Since $[L]$ is invertible in $\C \otimes_{\Z} \Z[G]$,
so are $x$ and $y$.
Since
\[
\C[G]^\times
= \left\{
\alpha\frac{[L^{\mathbf{b}}]}{[L^{\mathbf{a}}]}
\;\middle|\; \alpha \in \C^\times,\ 
\mathbf{a}, \mathbf{b} \in \Z_{\geq 0}^r
\right\},
\]
we have
\[
x = \alpha_1\frac{[L^{\mathbf{b}_1}]}{[L^{\mathbf{a}_1}]},
\qquad
y = \alpha_2\frac{[L^{\mathbf{b}_2}]}{[L^{\mathbf{a}_2}]}
\]
for $\alpha_1, \alpha_2 \in \C^\times$
and $\mathbf{a}_1, \mathbf{b}_1, \mathbf{a}_2, \mathbf{b}_2
\in \Z_{\geq 0}^r$.
Let $x = \sum_{i=1}^k \beta_i[M_i]$,
where $\beta_i \in \C^\times$, $M_i \in \Irr\mathscr{C}$, and
\[
\mathrm{PBW}(M_1) >_{\mathrm{lex}} \mathrm{PBW}(M_2)
>_{\mathrm{lex}} \cdots >_{\mathrm{lex}} \mathrm{PBW}(M_k).
\]
Then
\[
[L^{\mathbf{b}_1}]
= \sum_{i=1}^k \alpha_1^{-1}\beta_i
[L^{\mathbf{a}_1} \otimes M_i].
\]
By Lemma~\ref{Lem : tensor PBW},
$\alpha_1^{-1}\beta_1 = 1$ and
\[
\mathrm{PBW}(L^{\mathbf{b}_1})
= \mathrm{PBW}(L^{\mathbf{a}_1}) + \mathrm{PBW}(M_1).
\]
We prove $\alpha_1^{-1}\beta_s \in \Z$ by induction.
Assume $\alpha_1^{-1}\beta_1, \ldots,
\alpha_1^{-1}\beta_{s-1} \in \Z$.
Then
\[
[L^{\mathbf{b}_1}]
- \sum_{i=1}^{s-1}\alpha_1^{-1}\beta_i
[L^{\mathbf{a}_1} \otimes M_i]
= \sum_{i=s}^k \alpha_1^{-1}\beta_i
[L^{\mathbf{a}_1} \otimes M_i],
\]
and the left-hand side is in $K(\mathscr{C})_{\Z}$.
By Lemma~\ref{Lem : tensor PBW},
the coefficient of the simple module whose PBW parameter is
$\mathrm{PBW}(L^{\mathbf{a}_1}) + \mathrm{PBW}(M_s)$
in the right-hand side is $\alpha_1^{-1}\beta_s$.
Then we have $\alpha_1^{-1}\beta_s \in \Z$.
As a result, we have
\[
\alpha_1^{-1}x
= \frac{[L^{\mathbf{b}_1}]}{[L^{\mathbf{a}_1}]}
\in K(\mathscr{C})_{\Z}.
\]
By the same argument, we have
\[
\alpha_2^{-1}y
= \frac{[L^{\mathbf{b}_2}]}{[L^{\mathbf{a}_2}]}
\in K(\mathscr{C})_{\Z}.
\]
Then $\alpha_1^{-1}x = [X]$ and $\alpha_2^{-1}y = [Y]$
for some $X, Y \in \Irr\mathscr{C}$
by Proposition~\ref{Prop : divisibility by a real simple module}.
By the primeness of $L$, we conclude that
at least one of $X$ and $Y$ is trivial.
\end{proof}

\begin{Thm}
\label{Thm : Laurent families and maximal real sets}
Let $\{L_i \mid 1 \leq i \leq r\}$
be a real commuting family in $\mathscr{C}$.
Then it is a Laurent family if and only if it is a maximal real set.
\end{Thm}

\begin{proof}
Assume that $\{L_i\}$ is a Laurent family.
First, we prove that any $L_i$ is a prime module.
Let
\[
L_i \simeq \bigotimes_{k=1}^{\ell}M_k^{\otimes c_k}
\]
be a prime factorization.
Since $L_i$ commutes with any $L_j$ for $1 \leq j \leq r$,
so does each $M_k$.
Then, for each $1 \leq k \leq \ell$,
$M_k \simeq L^{\mathbf{a}_k}$
for some $\mathbf{a}_k \in \Z_{\geq 0}^r \setminus \{0\}$
by Definition~\ref{Def : Laurent family}(iii).
Then $L_i \simeq L^{\sum_{k=1}^{\ell}c_k\mathbf{a}_k}$.
By Definition~\ref{Def : Laurent family}(ii),
we have $\sum_{k=1}^{\ell}c_k\mathbf{a}_k = e_i$,
and we have that $L_i$ is prime.

Next, we prove the maximality.
Let $L \in \Irr\mathscr{C}$ be a real prime module
which commutes with any $L_i$.
Then $L \simeq L^{\mathbf{a}}$
for some $\mathbf{a} \in \Z_{\geq 0}^r$
by Definition~\ref{Def : Laurent family}(iii).
By the primeness of $L$, $L \simeq L_i$
for some $1 \leq i \leq r$.
Then we have that $\{L_i\}$ is a maximal real set.

If $\{L_i\}$ is a maximal real set, then it is a Laurent family
by Corollary~\ref{Cor : unique prime tensor factorization in C}
and Proposition~\ref{Prop : simple modules commuting with a maximal real set}.
\end{proof}

The following theorem is a monoidal analog of
\cite[Theorem 3.3]{BMRRT} and \cite[Theorem 2.4]{DK08}.

\begin{Thm}
\label{Thm : cardinality of maximal real sets}
Any maximal real set has exactly $N$ elements. 
Moreover, for a maximal real set $\{L_i \mid 1 \leq i \leq N\}$ and $t \in \mathbb{T}$, 
\[
\{\deg^t\iota([L_i]) \mid 1 \leq i \leq N\} \subset \Z^N
\]
is a $\Z$-basis of $\Z^N$.

\end{Thm}

\begin{proof}
It is immediate from Proposition~\ref{Prop : cardinality of Laurent families}, 
Propositoin~\ref{Prop : degrees of Laurent families form an integral basis}, 
and Theorem~\ref{Thm : Laurent families and maximal real sets}. 
\end{proof}

%%%%%%%%%%%%%%%%%%%%%%%%%%%%%%%%%%%%%%%%%%%%%%%%%%%%%%%%%%%%%%%%%%%%%%%%%%%%%%%%%%%%%%%%%%%%%%%%%%%%%%%%
\subsection{Exchange pair of maximal real sets}
\label{ssec : exchange pairs of maximal real sets}

\begin{Def}
\label{Def : almost complete maximal real sets}
A commuting family $\overline{\mathcal M}$ of real prime modules in
$\mathscr{C}$ is called an \emph{almost complete maximal real set} if
$|\overline{\mathcal M}| = N - 1$.
\end{Def}

\begin{Lem}
\label{Lem : exchange polynomial for maximal real sets}
Let $\overline{\mathcal M} = \{L_i \mid 2 \leq i \leq N\}$ be an almost
complete maximal real set.
Suppose that there exist real prime modules $L, L'$ such that
$L \not\simeq L'$ and $\{L\} \cup \overline{\mathcal M}$,
$\{L'\} \cup \overline{\mathcal M}$ are both maximal real sets.
Then there exists
$P(u_2, \ldots, u_N) \in \Z_{\geq 0}[u_2, \ldots, u_N]$ such that
\[
[L][L'] = P([L_2], \ldots, [L_N]).
\]
\end{Lem}

\begin{proof}
Since $\{L\} \cup \overline{\mathcal M}$ is a Laurent family,
$[L']$ has a Laurent expansion in
\[
\{[L], [L_2], \ldots, [L_N]\}.
\]
By Proposition~\ref{Prop : denominators and commutativity for Laurent families},
we have
\[
[L'] = \frac{Q([L], [L_2], \ldots, [L_N])}{[L]^k}
\]
for $k \in \Z_{\geq 0}$ and $Q \in \Z_{\geq 0}[u_1, \ldots, u_N]$
which is indivisible by $u_1$.
First we prove that $k = 1$.
If $k = 0$, then $L'$ commutes with $L$ by
Proposition~\ref{Prop : denominators and commutativity for Laurent families}.
Since $L'$ commutes with $L_2, \ldots, L_N$ by the assumption,
$L'$ commutes with $\{L\} \cup \overline{\mathcal M}$.
The maximality implies that $L'$ is a member of
$\{L\} \cup \overline{\mathcal M}$ and it is a contradiction.
Hence we have $k \geq 1$.
Consider the extension degree
\[
\bigl[\Q([L_i] \mid 2 \leq i \leq N)([L]) :
\Q([L_i] \mid 2 \leq i \leq N)([L'])\bigr].
\]
Since $[L']$ is a rational function in $\{[L], [L_2], \ldots, [L_N]\}$
whose denominator has degree $k$ as a polynomial in $[L]$,
the extension degree is at least $k$.
However, the extension degree is $1$ because both fields are
$\operatorname{Frac}(K(\mathscr{C})_{\Z})$ by
Proposition~\ref{Prop : Laurent expansions with respect to a Laurent family}.
Then we have $k = 1$.
By the same argument, we have
\[
[L] = \frac{Q'([L'], [L_2], \ldots, [L_N])}{[L']}
\]
for $Q' \in \Z_{\geq 0}[u_1, u_2, \ldots, u_N]$
which is indivisible by $u_1$.
Then
\[
Q([L], [L_2], \ldots, [L_N])
= Q'([L'], [L_2], \ldots, [L_N])
\]
holds.
By Corollary~\ref{Cor : unique prime tensor factorization in C},
it follows that $Q, Q' \in \Z_{\geq 0}[u_2, \ldots, u_N]$.
\end{proof}

\begin{Prop}
\label{Prop : at most two complements of an almost complete maximal real set}
Let $\overline{\mathcal M} = \{L_2, \ldots, L_N\}$ be an almost complete
maximal real set.
Then there are at most two real prime modules outside
$\overline{\mathcal M}$ which commute with every member of
$\overline{\mathcal M}$.
\end{Prop}

\begin{proof}
Suppose that there exist mutually non-isomorphic real prime modules
$L, L', L''$ which commute with every $L_i \in \overline{\mathcal M}$.

By Lemma~\ref{Lem : exchange polynomial for maximal real sets},
there exist $P', P'' \in \Z_{\geq 0}[u_2, \ldots, u_N]$ such that
\[
[L][L'] = P'([L_i] \mid 2 \leq i \leq N)
\quad\text{and}\quad
[L][L''] = P''([L_i] \mid 2 \leq i \leq N).
\]
Then we have
\[
[L']P''([L_i] \mid 2 \leq i \leq N)
= [L'']P'([L_i] \mid 2 \leq i \leq N).
\]
By the assumption, $L'$ is not isomorphic to $L''$ or $L_i$
for $2 \leq i \leq N$.
By Corollary~\ref{Cor : unique prime tensor factorization in C},
it is a contradiction.
\end{proof}

\begin{Def}
\label{Def : exchange pairs of maximal real sets}
Let $L, L' \in \Irr\mathscr{C}$ be non-isomorphic real prime modules.
We call $(L, L')$ an \emph{exchange pair} if there is an almost complete
maximal real set $\overline{\mathcal M}$ such that
$\{L\} \cup \overline{\mathcal M}$ and
$\{L'\} \cup \overline{\mathcal M}$ are both maximal real sets.
\end{Def}

\begin{Lem}
\label{Lem : unimodularity of PBW and tilde Lambda matrices}
Let $\mathcal M = \{L_i \mid 1 \leq i \leq N\}$ be a maximal real set.
Let
\[
P_{\mathcal M} \coloneq
(\mathrm{PBW}(L_i))_{1 \leq i \leq N} \in M_N(\Z).
\]
\begin{enumerate}
\item[(i)] $P_{\mathcal M} \in \mathrm{GL}_N(\Z)$.
\item[(ii)] $\widetilde{\Lambda}_{\mathcal M} \in \mathrm{GL}_N(\Z)$.
\end{enumerate}
\end{Lem}

\begin{proof}
(i) We can prove this similarly by replacing $\deg^{\mathcal M_t}$
by $\mathrm{PBW}$ and $M_{i;t}$ by $S_i$ in the proof of
Proposition~\ref{Prop : degrees of Laurent families form an integral basis}.

(ii) By Lemma~\ref{Lem : Lambda and PBW}, we have
\[
\Lambda_{\mathcal M}
= P_{\mathcal M}^{T}(\lambda_{kl})_{1 \leq k, l \leq N}P_{\mathcal M}.
\]
By Lemma~\ref{Lem : Lambda infinity tensor products and subquotients},
\[
\begin{aligned}
\Lambda^{\infty}(L_i, L_j)
&= \Lambda^{\infty}\left(
\bigotimes_{k=1}^{N} S_k^{\otimes\mathrm{PBW}(L_i)_k},
\bigotimes_{k=1}^{N} S_k^{\otimes\mathrm{PBW}(L_j)_k}
\right) \\
&= \mathrm{PBW}(L_i)^{T}
(\Lambda^{\infty}(S_k, S_l))_{1 \leq k, l \leq N}
\mathrm{PBW}(L_j).
\end{aligned}
\]
We have
\[
\Lambda^{\infty}_{\mathcal M}
= P_{\mathcal M}^{T}(- (\beta_k, \beta_l))_{1 \leq k, l \leq N}
P_{\mathcal M}
\]
by Lemma~\ref{Lem : Lambda infinity of cuspidal modules}.
Then we obtain that
\[
\widetilde{\Lambda}_{\mathcal M}
= P_{\mathcal M}^{T}
(\delta_{kl} + \delta(k < l)(\beta_k, \beta_l))_{1 \leq k, l \leq N}
P_{\mathcal M}.
\]
Since $(\delta_{kl} + \delta(k < l)(\beta_k, \beta_l))_{1 \leq k, l \leq N}$
is unitriangular, it is in $\mathrm{GL}_N(\Z)$.
Then $\widetilde{\Lambda}_{\mathcal M} \in \mathrm{GL}_N(\Z)$ by (i).
\end{proof}

\begin{Thm}
\label{Thm : exact sequence for an exchange pair}
Let $(L, L')$ be an exchange pair associated to an almost complete
maximal real set
$\overline{\mathcal{M}} = \{L_i \mid 2 \leq i \leq N\}$.
Let $\mathcal{M} \coloneq \{L_1 \coloneq L, L_2, \ldots, L_N\}$.
\begin{enumerate}
    \item[(i)] $\mathfrak{d}(L, L') = 1$.
    \item[(ii)] Let
    $\mathbf{v}_1 \coloneq -\widetilde{\Lambda}_{\mathcal{M}}^{-1}e_1$.
    Then there is an exact sequence
    \[
    0 \longrightarrow L^{[\mathbf{v}_1]_+}
    \longrightarrow L \otimes L'
    \longrightarrow L^{[-\mathbf{v}_1]_+}
    \longrightarrow 0.
    \]
\end{enumerate}
\end{Thm}

\begin{proof}
By Lemma~\ref{Lem : exchange polynomial for maximal real sets},
\[
[L \otimes L'] = \sum_{s=1}^{\ell}c_s[L^{\mathbf{b}(s)}]
\]
for $c_s \in \Z_{>0}$, $\mathbf{b}(s) \in \Z_{\geq 0}^N$,
and $\mathbf{b}(s)_1 = 0$.
Let $L \nabla L' \simeq L^{\mathbf{h}}$ and
$L \Delta L' \simeq L^{\mathbf{s}}$.
Note that the coefficients of $[L^{\mathbf{h}}]$ and $[L^{\mathbf{s}}]$
are $1$ by Lemma~\ref{Lem : Lambda inequalities for heads and socles}.
By Proposition~\ref{Prop : d vanishing and commutativity}
and Lemma~\ref{Lem : Lambda inequalities for heads and socles},
\[
\begin{aligned}
\bigl(\Lambda_{\mathcal{M}}(\mathbf{s} - \mathbf{h})\bigr)_1
&= \Lambda(L, L \Delta L') - \Lambda(L, L \nabla L') \\
&= -\Lambda(L \Delta L', L) - \Lambda(L, L') \\
&= -\Lambda(L', L) - \Lambda(L, L') \\
&= -2\mathfrak{d}(L, L').
\end{aligned}
\]
Since $L_i$ commutes with $L$ and $L'$ for any $2 \leq i \leq N$,
\[
\begin{aligned}
\bigl(\Lambda_{\mathcal{M}}(\mathbf{s} - \mathbf{h})\bigr)_i
&= \Lambda(L_i, L \Delta L') - \Lambda(L_i, L \nabla L') \\
&= \bigl(\Lambda(L_i, L) + \Lambda(L_i, L')\bigr)
 - \bigl(\Lambda(L_i, L) + \Lambda(L_i, L')\bigr) \\
&= 0
\end{aligned}
\]
by Proposition~\ref{Prop : commutativity with simple subquotients}.
By Lemma~\ref{Lem : Lambda infinity tensor products and subquotients},
$\Lambda^{\infty}_{\mathcal{M}}(\mathbf{s} - \mathbf{h}) = 0$.
Hence
\[
\widetilde{\Lambda}_{\mathcal{M}}(\mathbf{s} - \mathbf{h})
= -\mathfrak{d}(L, L')e_1.
\]
By Proposition~\ref{Prop : Lambda inequalities for Laurent expansions},
for $1 \leq s \leq \ell$ such that
$\mathbf{b}(s) \neq \mathbf{s}, \mathbf{h}$,
there exists $a_s \in \Z$ with $0 < a_s < \mathfrak{d}(L, L')$
such that
\[
\widetilde{\Lambda}_{\mathcal{M}}(\mathbf{b}(s) - \mathbf{h})
= -a_s e_1.
\]
Multiplying by $\widetilde{\Lambda}_{\mathcal{M}}^{-1}$, we have
\[
\mathbf{s} - \mathbf{h} = \mathfrak{d}(L, L')\mathbf{v}_1,
\qquad
\mathbf{b}(s) - \mathbf{h} = a_s\mathbf{v}_1.
\]
Note that $(\mathbf{v}_1)_1 = 0$
because $(\mathbf{s} - \mathbf{h})_1 = 0$.
Let
\[
P(t) \coloneq 1
+ \sum_{\mathbf{b}(s) \neq \mathbf{s}, \mathbf{h}}c_s t^{a_s}
+ t^{\mathfrak{d}(L, L')}
\in \Z_{\geq 0}[t].
\]
Then
\[
[L \otimes L'] = [L^{\mathbf{h}}]
P\left(\prod_{i=1}^N[L_i]^{(\mathbf{v}_1)_i}\right)
\]
by the above argument.
Let
\[
P(t) = \prod_{j=1}^{\mathfrak{d}(L, L')}(t - \alpha_j)
\qquad \text{in $\C[t]$}.
\]
Then
\[
[L][L'] = [L^{\mathbf{h}}]
\prod_{j=1}^{\mathfrak{d}(L, L')}
\left(\prod_{i=1}^N[L_i]^{(\mathbf{v}_1)_i} - \alpha_j\right)
\]
in $\C \otimes K(\mathscr{C})_{\Z}
[[L_i]^{-1} \mid 2 \leq i \leq N]$.

(i) Since $(L, L')$ is an exchange pair,
$\mathfrak{d}(L, L') > 0$ is clear.
Since $[L']$ and $[L_i]$ for $1 \leq i \leq N$ are prime elements
in $\C \otimes K(\mathscr{C})_{\Z}$ and
$[L']$ is not associate to $[L_i]$ for any $1 \leq i \leq N$
by Lemma~\ref{Lem : real prime modules give prime elements over C},
$[L']$ is a prime element in
$\C \otimes K(\mathscr{C})_{\Z}
[[L_i]^{-1} \mid 1 \leq i \leq N]$.
In $\C \otimes K(\mathscr{C})_{\Z}
[[L_i]^{-1} \mid 1 \leq i \leq N]$, we have
\[
[L'] = \frac{[L^{\mathbf{h}}]}{[L_1]}
\prod_{j=1}^{\mathfrak{d}(L, L')}
\left(\prod_{i=1}^N[L_i]^{(\mathbf{v}_1)_i} - \alpha_j\right).
\]
Since $\mathbf{v}_1 \neq 0$ and $\alpha_j \in \C^{\times}$,
$\prod_{i=1}^N[L_i]^{(\mathbf{v}_1)_i} - \alpha_j$
is not invertible in
$\C \otimes K(\mathscr{C})_{\Z}
[[L_i]^{-1} \mid 1 \leq i \leq N]$.
Then the right-hand side is a product of at least
$\mathfrak{d}(L, L')$ prime elements.
Therefore $\mathfrak{d}(L, L') \leq 1$.

(ii) We have
\[
[L \otimes L'] = [L^{\mathbf{h}}]
\left(1 + \prod_{i=1}^N[L_i]^{(\mathbf{v}_1)_i}\right)
\]
by (i).
Since $\mathbf{h} + \mathbf{v}_1 = \mathbf{s}$,
$\mathbf{h}, \mathbf{h} + \mathbf{v}_1 \in \Z_{\geq 0}^N$.
Note that $h_1 = 0$
by Lemma~\ref{Lem : exchange polynomial for maximal real sets}.
Suppose that $\min\{h_i, (\mathbf{v}_1 + \mathbf{h})_i\} > 0$
for some $2 \leq i \leq N$.
Then $[L_i] \mid [L][L']$ in $K(\mathscr{C})_{\Z}$.
However, $[L_i]$, $[L]$, and $[L']$ are non-associate prime elements
in $K(\mathscr{C})_{\Z}$, which is a contradiction.
Therefore $\min\{h_i, (\mathbf{v}_1 + \mathbf{h})_i\} = 0$
for any $1 \leq i \leq N$.
Thus we have $\mathbf{h} = [-\mathbf{v}_1]_+$ and
$\mathbf{v}_1 + \mathbf{h} = [\mathbf{v}_1]_+$.
Then we obtain the exact sequence
\[
0 \longrightarrow L^{[\mathbf{v}_1]_+}
\longrightarrow L \otimes L'
\longrightarrow L^{[-\mathbf{v}_1]_+}
\longrightarrow 0.
\]
\end{proof}

The following theorem is a monoidal analog of
\cite[Theorem 7.5]{BMRRT}, \cite[Theorem 5.8]{FG}.

\begin{Thm}
\label{Thm : characterization of exchange pairs}
Let $L, L' \in \Irr\mathscr{C}$ be non-isomorphic real prime modules.
Then $(L, L')$ is an exchange pair if and only if
$\mathfrak{d}(L, L') = 1$.
\end{Thm}

\begin{proof}
The only if implication is shown in
Theorem~\ref{Thm : exact sequence for an exchange pair}.
Assume that $\mathfrak{d}(L, L') = 1$.
By Lemma~\ref{Lem : d one exact sequence},
\[
0 \longrightarrow L \Delta L'
\longrightarrow L \otimes L'
\longrightarrow L \nabla L'
\longrightarrow 0
\]
is exact and $L \Delta L'$ and $L \nabla L'$ are both real.
We have
\[
\mathfrak{d}(L, L \Delta L')
= \mathfrak{d}(L, L \nabla L')
= \mathfrak{d}(L', L \Delta L')
= \mathfrak{d}(L', L \nabla L') = 0
\]
by Lemma~\ref{Lem : decrease of d}.
In addition, we have
$\mathfrak{d}(L \Delta L', L \nabla L') = 0$
by Lemma~\ref{Lem : Lambda tensor products and subquotients}.
Thus $L \otimes (L \Delta L') \otimes (L \nabla L')$ is a real module. 
Let 
\[
\mathcal{S} \coloneq
\{\text{prime factors of }
L \otimes (L \Delta L') \otimes (L \nabla L')\}
\]
be a commuting family of real prime modules in $\mathscr{C}$.
Let $\mathcal{M} = \{L_1 \coloneq L, L_2, \ldots, L_N\}$
be a maximal real set which contains $\mathcal{S}$.
Since $L \Delta L'$ and $L \nabla L'$ are monomials in $\mathcal{M}$,
\[
[L'] = \frac{[L^{\mathbf{a}}] + [L^{\mathbf{b}}]}{[L]}
\qquad \text{for $\mathbf{a}, \mathbf{b} \in \Z_{\geq 0}^N$}.
\]
By Proposition~\ref{Prop : denominators and commutativity for Laurent families},
$L'$ commutes with any $L_i$ for $2 \leq i \leq N$.
Thus $\{L'\} \cup \{L_i \mid 2 \leq i \leq N\}$
is a maximal real set.
Then we have that $(L, L')$ is an exchange pair.
\end{proof}

\subsection{$\Lambda$-admissible structure of maximal real sets}
\label{ssec : Lambda admissible structures of maximal real sets}

Let $J \coloneq \{1, \ldots, N\}$, and let
$J_{\mathrm{uf}}$ and $J_{\mathrm{fr}}$ be as in
Subsection~\ref{ssec : Lambda cluster algebra structure}.
By Proposition~\ref{Prop : frozen commutativity} and
Proposition~\ref{Prop : cluster monomials and real prime modules},
any maximal real set contains the real prime modules corresponding to
the frozen variables.
In what follows, by changing the labeling if necessary, we may assume
that each maximal real set $\mathcal M = \{L_i \mid i \in J\}$ satisfies
$L_j \simeq M_{j;t_0}$ for all $j \in J_{\mathrm{fr}}$.

For a maximal real set $\mathcal M = \{L_i \mid i \in J\}$, let
\[
\widetilde B_{\mathcal M}
\coloneq -\widetilde\Lambda_{\mathcal M}^{-1}
\begin{pmatrix} I_{|J_{\mathrm{uf}}|} \\ 0 \end{pmatrix}.
\]
Here,
\[
\begin{pmatrix} I_{|J_{\mathrm{uf}}|} \\ 0 \end{pmatrix}
= (\delta_{ij})_{i \in J,\,j \in J_{\mathrm{uf}}}.
\]

\begin{Prop}
\label{Prop : exchange matrices of monoidal clusters}
Let $t \in \mathbb T$, and let
$\mathcal M_t = \{M_{i;t} \mid i \in J\}$ be the corresponding
monoidal cluster.
Then $\widetilde B_t = \widetilde B_{\mathcal M_t}$.
\end{Prop}

\begin{proof}
Let $k \in J_{\mathrm{uf}}$.
Since $(M_{k;t}, M_{k;\mu_k t})$ is an exchange pair corresponding to
$\{M_{i;t} \mid i \in J \setminus \{k\}\}$,
\[
0 \longrightarrow
M_t^{[\widetilde B_{\mathcal M_t}e_k]_+}
\longrightarrow M_{k;t} \otimes M_{k;\mu_k t}
\longrightarrow M_t^{[-\widetilde B_{\mathcal M_t}e_k]_+}
\longrightarrow 0
\]
is an exact sequence by
Theorem~\ref{Thm : exact sequence for an exchange pair}.
By Definition~\ref{Def : Lambda admissible monoidal seed},
\[
0 \longrightarrow M_t^{[\widetilde B_t e_k]_+}
\longrightarrow M_{k;t} \otimes M_{k;\mu_k t}
\longrightarrow M_t^{[-\widetilde B_t e_k]_+}
\longrightarrow 0
\]
is also an exact sequence.
By Proposition~\ref{Prop : simplicity of heads and socles},
these exact sequences are non-split.
Then we see that
\[
M_t^{[\widetilde B_{\mathcal M_t}e_k]_+}
\simeq M_t^{[\widetilde B_t e_k]_+},
\qquad
M_t^{[-\widetilde B_{\mathcal M_t}e_k]_+}
\simeq M_t^{[-\widetilde B_t e_k]_+}.
\]
By Corollary~\ref{Cor : unique prime tensor factorization in C},
$\widetilde B_{\mathcal M_t}e_k = \widetilde B_t e_k$ holds.
\end{proof}

\begin{Prop}
\label{Prop : compatibility for maximal real sets}
Let $\mathcal M = \{L_i \mid i \in J\}$ be a maximal real set.
\begin{enumerate}
    \item[(i)] $\widetilde B_{\mathcal M}$ is an exchange matrix.
    \item[(ii)] $(\widetilde B_{\mathcal M}, \Lambda_{\mathcal M})$
    is a compatible pair of type $2I_{|J_{\mathrm{uf}}|}$.
    \item[(iii)]
    $\Lambda^{\infty}_{\mathcal M}\widetilde B_{\mathcal M} = 0$.
\end{enumerate}
\end{Prop}

\begin{proof}
Let $P_{\mathcal M} = (\mathrm{PBW}(L_i))_{i \in J}$ and
$P_{t_0} = (\mathrm{PBW}(M_{i;t_0}))_{i \in J}$.
Note that the claim is clear if $\mathcal M$ is a monoidal cluster.

(i) By the proof of
Lemma~\ref{Lem : unimodularity of PBW and tilde Lambda matrices},
we have
\[
\widetilde\Lambda_{\mathcal M}
= P_{\mathcal M}^{T}(P_{t_0}^{T})^{-1}
\widetilde\Lambda_{t_0}P_{t_0}^{-1}P_{\mathcal M}.
\]
Thus
\[
\begin{aligned}
\left.\widetilde B_{\mathcal M}\right|_
{J_{\mathrm{uf}} \times J_{\mathrm{uf}}}
&= (I_{|J_{\mathrm{uf}}|}\ 0)\widetilde B_{\mathcal M} \\
&= -(I_{|J_{\mathrm{uf}}|}\ 0)\widetilde\Lambda_{\mathcal M}^{-1}
\begin{pmatrix} I_{|J_{\mathrm{uf}}|} \\ 0 \end{pmatrix} \\
&= -\bigl((I_{|J_{\mathrm{uf}}|}\ 0)P_{\mathcal M}^{-1}P_{t_0}\bigr)
\widetilde\Lambda_{t_0}^{-1}
\bigl((I_{|J_{\mathrm{uf}}|}\ 0)P_{\mathcal M}^{-1}P_{t_0}\bigr)^{T}.
\end{aligned}
\]
Since $L_j \simeq M_{j;t_0}$ for $j \in J_{\mathrm{fr}}$,
$P_{t_0}e_j = P_{\mathcal M}e_j$ for $j \in J_{\mathrm{fr}}$.
Therefore
\[
\left.(P_{\mathcal M}^{-1}P_{t_0})\right|_
{J \times J_{\mathrm{fr}}}
= \begin{pmatrix} 0 \\ I_{|J_{\mathrm{fr}}|} \end{pmatrix}.
\]
Since $P_{\mathcal M}^{-1}P_{t_0} \in \mathrm{GL}_N(\Z)$, we have
\[
A_{\mathrm{uf}} \coloneq
\left.(P_{\mathcal M}^{-1}P_{t_0})\right|_
{J_{\mathrm{uf}} \times J_{\mathrm{uf}}}
\in \mathrm{GL}_{|J_{\mathrm{uf}}|}(\Z).
\]
Then
\[
(I_{|J_{\mathrm{uf}}|}\ 0)(P_{\mathcal M}^{-1}P_{t_0})
= (A_{\mathrm{uf}}\ 0)
= A_{\mathrm{uf}}(I_{|J_{\mathrm{uf}}|}\ 0).
\]
Hence it holds that
\[
\begin{aligned}
\left.\widetilde B_{\mathcal M}\right|_
{J_{\mathrm{uf}} \times J_{\mathrm{uf}}}
&= -\bigl(A_{\mathrm{uf}}(I_{|J_{\mathrm{uf}}|}\ 0)\bigr)
\widetilde\Lambda_{t_0}^{-1}
\begin{pmatrix} I_{|J_{\mathrm{uf}}|} \\ 0 \end{pmatrix}
A_{\mathrm{uf}}^{T} \\
&= A_{\mathrm{uf}}(I_{|J_{\mathrm{uf}}|}\ 0)
\widetilde B_{t_0}A_{\mathrm{uf}}^{T}.
\end{aligned}
\]
Because $\widetilde B_{t_0}$ is an exchange matrix,
$(I_{|J_{\mathrm{uf}}|}\ 0)\widetilde B_{t_0}$ is skew-symmetric,
and thus
$\left.\widetilde B_{\mathcal M}\right|_
{J_{\mathrm{uf}} \times J_{\mathrm{uf}}}$ is skew-symmetric.

(ii) By the proof of
Lemma~\ref{Lem : unimodularity of PBW and tilde Lambda matrices},
we have
\[
\Lambda_{\mathcal M}
= P_{\mathcal M}^{T}(P_{t_0}^{T})^{-1}
\Lambda_{t_0}P_{t_0}^{-1}P_{\mathcal M}.
\]
Thus it holds that
\[
\begin{aligned}
\widetilde B_{\mathcal M}^{T}\Lambda_{\mathcal M}
&= -(I_{|J_{\mathrm{uf}}|}\ 0)
(\widetilde\Lambda_{\mathcal M}^{T})^{-1}\Lambda_{\mathcal M}
\\
&= -(I_{|J_{\mathrm{uf}}|}\ 0)(P_{\mathcal M}^{-1}P_{t_0})
(\widetilde\Lambda_{t_0}^{T})^{-1}
(P_{\mathcal M}^{-1}P_{t_0})^{T} 
 (P_{t_0}^{-1}P_{\mathcal M})^{T}
\Lambda_{t_0}(P_{t_0}^{-1}P_{\mathcal M}) \\
&= -(I_{|J_{\mathrm{uf}}|}\ 0)(P_{\mathcal M}^{-1}P_{t_0})
(\widetilde\Lambda_{t_0}^{T})^{-1}
\Lambda_{t_0}(P_{t_0}^{-1}P_{\mathcal M}).
\end{aligned}
\]
Note that $P_{\mathcal M}^{-1}P_{t_0}$ is of the form
\[
\begin{pmatrix} A_{\mathrm{uf}} & 0 \\ * & I_{|J_{\mathrm{fr}}|} \end{pmatrix}.
\]
Hence we have
\[
\begin{aligned}
&-(I_{|J_{\mathrm{uf}}|}\ 0)(P_{\mathcal M}^{-1}P_{t_0})
(\widetilde\Lambda_{t_0}^{T})^{-1}
\Lambda_{t_0}(P_{t_0}^{-1}P_{\mathcal M}) \\
&\quad= -A_{\mathrm{uf}}(I_{|J_{\mathrm{uf}}|}\ 0)
(\widetilde\Lambda_{t_0}^{T})^{-1}\Lambda_{t_0}
\begin{pmatrix}
A_{\mathrm{uf}}^{-1} & 0 \\ * & I_{|J_{\mathrm{fr}}|}
\end{pmatrix} \\
&\quad= A_{\mathrm{uf}}\widetilde B_{t_0}^{T}\Lambda_{t_0}
\begin{pmatrix}
A_{\mathrm{uf}}^{-1} & 0 \\ * & I_{|J_{\mathrm{fr}}|}
\end{pmatrix} \\
&\quad= 2(A_{\mathrm{uf}}\ 0)
\begin{pmatrix}
A_{\mathrm{uf}}^{-1} & 0 \\ * & I_{|J_{\mathrm{fr}}|}
\end{pmatrix} \\
&\quad= (2I_{|J_{\mathrm{uf}}|}\ 0).
\end{aligned}
\]

(iii) By the proof of
Lemma~\ref{Lem : unimodularity of PBW and tilde Lambda matrices},
we have
\[
\Lambda^{\infty}_{\mathcal M}
= (P_{t_0}^{-1}P_{\mathcal M})^{T}\Lambda^{\infty}_{t_0}
(P_{t_0}^{-1}P_{\mathcal M}).
\]
The proof is the same as (ii), except using the fact that
$\Lambda^{\infty}_{t_0}\widetilde B_{t_0} = 0$.
\end{proof}

The goal of this subsection is to prove that
$(\mathcal M, \widetilde B_{\mathcal M})$ is a completely
$\Lambda$-admissible monoidal seed for all maximal real sets
$\mathcal M = \{L_i \mid i \in J\}$.
For $k \in J_{\mathrm{uf}}$, let
$\mathbf v_k \coloneq \widetilde B_{\mathcal M}e_k$.
What remains is to prove the existence of a real prime module $L'_k$
such that
\[
0 \longrightarrow L^{[\mathbf v_k]_+}
\longrightarrow L_k \otimes L'_k
\longrightarrow L^{[-\mathbf v_k]_+}
\longrightarrow 0
\]
is exact.
By Proposition~\ref{Prop : compatibility for maximal real sets},
$(\mathbf v_k)_k = 0$ and
\[
\Lambda(L_k, L^{[\mathbf v_k]_+})
- \Lambda(L_k, L^{[-\mathbf v_k]_+})
= e_k^{T}\Lambda_{\mathcal M}\widetilde B_{\mathcal M}e_k
= -2.
\]

\begin{Lem}
\label{Lem : binomial divisibility by a real simple module}
Let $L \in \Irr\mathscr{C}$ be a real module, and let
$S, H \in \Irr\mathscr{C}$ such that
\[
\Lambda(L, S) - \Lambda(L, H) = -2,
\qquad
\mathfrak{d}(L, S) = \mathfrak{d}(L, H) = 0.
\]
Assume that $[L] \nmid [S], [H]$ in
$\C \otimes_{\Z} K(\mathscr{C})_{\Z}$.
Suppose that there exist
$\phi \in \C \otimes_{\Z} K(\mathscr{C})_{\Z}$
and $\alpha \in \C^{\times}$ such that
\[
[L]\phi = [S] - \alpha[H]
\]
in $\C \otimes_{\Z} K(\mathscr{C})_{\Z}$.
Then we have $\alpha = -1$, and there exists
$L' \in \Irr\mathscr{C}$ such that $\phi = [L']$.
Moreover, $\mathfrak{d}(L, L') = 1$, and it holds that
\[
0 \longrightarrow S \longrightarrow L \otimes L'
\longrightarrow H \longrightarrow 0
\]
is an exact sequence.
\end{Lem}

\begin{proof}
We prove it in a similar way to
\cite[Theorem 4.1]{KK19} and \cite[Theorem 4.13]{KKOP20}.
Let
\[
\phi = \sum_{s \in K} c_s[X_s]
\]
for $c_s \in \C^{\times}$ and $X_s \in \Irr\mathscr{C}$.
It follows that
\[
\sum_{s \in K} c_s[L \otimes X_s] = [S] - \alpha[H].
\]
Take $s_0 \in K$ such that
\[
\Lambda(X_{s_0}, L)
= \max\{\Lambda(X_s, L) \mid s \in K\}.
\]
Suppose that there exists $s_0' \in K \setminus\{s_0\}$
such that $L \Delta X_{s_0}$ is a simple subquotient of
$L \otimes X_{s_0'}$.
Note that $L \Delta X_{s_0'} \not\simeq L \Delta X_{s_0}$
by Lemma~\ref{Lem : hd soc injectivity}.
By Lemma~\ref{Lem : Lambda inequalities for heads and socles},
it follows that
\[
\begin{aligned}
\Lambda(X_{s_0}, L)
&= \Lambda(L \Delta X_{s_0}, L) \\
&< \Lambda(L \Delta X_{s_0'}, L) \\
&= \Lambda(X_{s_0'}, L).
\end{aligned}
\]
It is a contradiction to the choice of $s_0$.
Thus we have $L \Delta X_{s_0} \simeq S$ or $H$,
and the coefficient of $[L \Delta X_{s_0}]$
in the right-hand side is $c_{s_0}$.
Since
\[
\begin{aligned}
&\max\{\Lambda(T, L) \mid
T \text{ is a simple subquotient of } X_s \otimes L
\text{ for some } s \in K\} \\
&\qquad = \Lambda(L \Delta X_{s_0}, L)
\end{aligned}
\]
by Lemma~\ref{Lem : Lambda inequalities for heads and socles}
and $\Lambda(S, L) = \Lambda(H, L) + 2$ by the assumption,
$L \Delta X_{s_0} \simeq S$ holds.
Thus we have $c_{s_0} = 1$.
Since $[L] \nmid [S]$ by the assumption,
we have $\mathfrak{d}(L, X_{s_0}) \geq 1$.
Take $t_0 \in K$ such that
\[
\Lambda(L, X_{t_0})
= \max\{\Lambda(L, X_s) \mid s \in K\}.
\]
By the same argument, we have $L \nabla X_{t_0} \simeq H$
and $c_{t_0} = -\alpha$.
Suppose that there exists $t_1 \neq t_0$ such that
$\Lambda(L, X_{t_1}) = \Lambda(L, X_{t_0})$.
Then $L \nabla X_{t_0} \simeq L \nabla X_{t_1} \simeq H$
holds, and it is a contradiction to
Lemma~\ref{Lem : hd soc injectivity}.

Let us prove that $s_0 = t_0$.
It is enough to prove
$\Lambda(L, X_{s_0}) = \Lambda(L, X_{t_0})$.
By Lemma~\ref{Lem : Lambda inequalities for heads and socles}
and the assumption, we have
\[
\begin{aligned}
\Lambda(L, X_{s_0})
&= 2\mathfrak{d}(L, X_{s_0}) - \Lambda(X_{s_0}, L) \\
&= 2\mathfrak{d}(L, X_{s_0})
   - \Lambda(L \Delta X_{s_0}, L) \\
&= 2\mathfrak{d}(L, X_{s_0}) - \Lambda(S, L) \\
&= 2\mathfrak{d}(L, X_{s_0}) + \Lambda(L, S) \\
&\geq 2 + \Lambda(L, S) \\
&= \Lambda(L, H), \\
\Lambda(L, X_{t_0})
&= \Lambda(L, L \nabla X_{t_0}) = \Lambda(L, H).
\end{aligned}
\]
By the choice of $t_0$, we have
$\Lambda(L, X_{s_0}) = \Lambda(L, X_{t_0})$.
Therefore, we have $-\alpha = c_{t_0} = c_{s_0} = 1$ and
\[
\begin{aligned}
\mathfrak{d}(L, X_{s_0})
&= \frac{1}{2}\bigl(
\Lambda(L, L \nabla X_{s_0})
+ \Lambda(L \Delta X_{s_0}, L)\bigr) \\
&= \frac{1}{2}\bigl(\Lambda(L, H) + \Lambda(S, L)\bigr) \\
&= \frac{1}{2}\bigl(\Lambda(L, H) - \Lambda(L, S)\bigr) \\
&= 1.
\end{aligned}
\]
By Lemma~\ref{Lem : d one exact sequence},
we have an exact sequence
\[
0 \longrightarrow S \longrightarrow L \otimes X_{s_0}
\longrightarrow H \longrightarrow 0.
\]
\end{proof}

\begin{Lem}
\label{Lem : Lambda admissibility and binomial divisibility}
Let $\mathcal{M} = \{L_i \mid i \in J\}$ be a maximal real set.
Then the following conditions are equivalent:
\begin{enumerate}
    \item[(i)] $(\mathcal{M}, \widetilde{B}_{\mathcal{M}})$
    is a $\Lambda$-admissible monoidal seed.
    \item[(ii)] For each $k \in J_{\mathrm{uf}}$,
    there exists $\alpha_k \in \C^{\times}$ such that
    \[
    \frac{[L^{[\mathbf{v}_k]_+}]
    - \alpha_k[L^{[-\mathbf{v}_k]_+}]}{[L_k]}
    \in \C \otimes K(\mathscr{C})_{\Z}.
    \]
\end{enumerate}
\end{Lem}

\begin{proof}
The implication (i) $\Rightarrow$ (ii) follows from
Definition~\ref{Def : Lambda admissible monoidal seed}.
Assume (ii).
By Lemma~\ref{Lem : binomial divisibility by a real simple module},
there exists $L'_k \in \Irr\mathscr{C}$ such that
$\mathfrak{d}(L_k, L'_k) = 1$ and
\[
0 \longrightarrow L^{[\mathbf{v}_k]_+}
\longrightarrow L_k \otimes L'_k
\longrightarrow L^{[-\mathbf{v}_k]_+}
\longrightarrow 0
\]
is exact.
By Proposition~\ref{Prop : denominators and commutativity for Laurent families},
$L'_k$ commutes with all $L_i$ for $i \neq k$.
Therefore it is enough to prove that $L'_k$ is real prime.

First we prove the primeness of $L'_k$.
Suppose $L'_k$ is not prime.
Since $\mathfrak{d}(L_k, L'_k) = 1$,
there exists a prime factor $P$ such that $\mathfrak{d}(L_k, P) = 0$.
Note that we can write $L'_k \simeq P \otimes P'$.
Since $\mathfrak{d}(L_k, P') = 1$,
the length of $L_k \otimes P'$ is at least $2$.
Because $L_k \otimes P \otimes P'$ has length $2$,
we have that the length of $L_k \otimes P'$ is two and
$P$ strongly commutes with both simple subquotients.
Thus we have
\[
[P] \mid [L^{[\mathbf{v}_k]_+}],\ [L^{[-\mathbf{v}_k]_+}].
\]
It contradicts Proposition~\ref{Prop : real prime modules give prime elements}.

Next we prove the reality of $L'_k$.
Suppose $L'_k$ is imaginary.
Since $(\mathbf{v}_k)_k = 0$,
$L'_k$ commutes with $L^{[\mathbf{v}_k]_+}$ and
$L^{[-\mathbf{v}_k]_+}$.
Then we have
\[
[L_k \otimes L'_k \otimes L'_k]
= [L^{[\mathbf{v}_k]_+} \otimes L'_k]
+ [L^{[-\mathbf{v}_k]_+} \otimes L'_k],
\]
and we obtain that $L_k \otimes L'_k \otimes L'_k$ has length $2$.
Because $L'_k$ is imaginary,
the length of $L'_k \otimes L'_k$ is at least $2$.
Since $L_k \otimes L'_k \otimes L'_k$ and $L'_k \otimes L'_k$
have the same length,
$L_k$ commutes with all simple subquotients of $L'_k \otimes L'_k$.
Then we obtain that
\[
[L_k] \mid [L^{[\mathbf{v}_k]_+} \otimes L'_k]
\qquad\text{in $K(\mathscr{C})_{\Z}$}.
\]
By Proposition~\ref{Prop : real prime modules give prime elements}, we have
\[
[L_k] \mid [L'_k]
\qquad\text{in $K(\mathscr{C})_{\Z}$}.
\]
By Proposition~\ref{Prop : divisibility by a real simple module},
there exists $Y \in \Irr\mathscr{C}$ such that
$L'_k \simeq L_k \otimes Y$.
It contradicts the primeness of $L'_k$.
\end{proof}

The following theorem is a monoidal analog of
\cite[Theorem 5.1]{BMRRT}, \cite[Theorem 5.3]{IY08}.

\begin{Thm}
\label{Thm : Lambda admissibility of maximal real sets}
Let $\mathcal{M} = \{L_i \mid i \in J\}$ be a maximal real set.
Then $(\mathcal{M}, \widetilde{B}_{\mathcal{M}})$ is a
$\Lambda$-admissible monoidal seed.
\end{Thm}

\begin{proof}
Let $k \in J_{\mathrm{uf}}$.
First we prove that there is a simple module which does not commute
with $L_k$.
Let $t \in \mathbb{T}$ and $\mathcal{M}_t$ be the corresponding
monoidal cluster.
If $L_k \notin \mathcal{M}_t$, then $L_k$ does not commute with
some element of $\mathcal{M}_t$ because $\mathcal{M}_t$ is a
maximal real set.
Suppose that there exists $i \in J$ such that $L_k \simeq M_{i;t}$.
Since $k \in J_{\mathrm{uf}}$, $L_k$ does not correspond to
frozen variables and we have $i \in J_{\mathrm{uf}}$.
Then $M_{i; \mu_i t}$ does not commute with $L_k$. 

Let $X$ be a simple module which does not commute with $L_k$.
By Proposition~\ref{Prop : denominators and commutativity for Laurent families},
$\mathbf{d}(X)_k > 0$ holds.
Let us consider
\[
\left(\bigotimes_{i \neq k}
L_i^{\otimes\max\{0, \mathbf{d}(X)_i\}}\right) \otimes X.
\]
Since each coefficient of Laurent polynomials corresponding to modules
is positive, there exists some simple subquotient $S$ of
\[
\left(\bigotimes_{i \neq k}
L_i^{\otimes\max\{0, \mathbf{d}(X)_i\}}\right) \otimes X
\]
such that $\mathbf{d}(S)_k > 0$ and $\mathbf{d}(S)_i \leq 0$
for all $i \neq k$.

Let
\[
[L_k^{\otimes\mathbf{d}(S)_k} \otimes S]
= \sum_{s=1}^{\ell}c_s[L^{\mathbf{b}(s)}]
\]
for $c_s \in \Z_{>0}$ and $\mathbf{b}(s) \in \Z_{\geq 0}^N$.
Let
\[
L_k^{\otimes\mathbf{d}(S)_k} \nabla S \simeq L^{\mathbf{h}},
\qquad
L_k^{\otimes\mathbf{d}(S)_k} \Delta S \simeq L^{\mathbf{s}}.
\]
By the same argument as
Theorem~\ref{Thm : exact sequence for an exchange pair}, we have
\[
\mathbf{s} - \mathbf{h}
= \mathfrak{d}(L_k, S)\mathbf{v}_k
\]
and
\[
\mathbf{b}(s) - \mathbf{h} = a_s\mathbf{v}_k
\qquad\text{for }0 < a_s < \mathfrak{d}(L_k, S)
\]
for all $\mathbf{b}(s) \neq \mathbf{s}, \mathbf{h}$.
Thus there exists
\[
P(t) = 1 + \sum_{j=1}^{\mathfrak{d}(L_k, S)-1}a_jt^j
+ t^{\mathfrak{d}(L_k, S)} \in \Z_{\geq 0}[t]
\]
such that
\[
[L_k^{\otimes\mathbf{d}(S)_k} \otimes S]
= [L^{\mathbf{h}}]
P\left(\prod_{i \in J}[L_i]^{(\mathbf{v}_k)_i}\right).
\]
Let
\[
P(t) = \prod_{j=1}^{\mathfrak{d}(L_k, S)}(t - \alpha_j)
\qquad\text{for }\alpha_j \in \C^{\times}.
\]
Then it follows that
\[
[L_k^{\otimes\mathbf{d}(S)_k} \otimes S]
= [L^{\mathbf{h}}]
\prod_{j=1}^{\mathfrak{d}(L_k, S)}
\left(\prod_{i \in J}[L_i]^{(\mathbf{v}_k)_i} - \alpha_j\right).
\]
Multiplying by $[L^{[-\mathbf{v}_k]_+}]^{\mathfrak{d}(L_k, S)}$,
we obtain that
\[
\begin{aligned}
&[L^{\mathfrak{d}(L_k, S)[-\mathbf{v}_k]_+}
\otimes L_k^{\otimes\mathbf{d}(S)_k} \otimes S] \\
&\qquad = [L^{\mathbf{h}}]
\prod_{j=1}^{\mathfrak{d}(L_k, S)}
\left([L^{[\mathbf{v}_k]_+}]
- \alpha_j[L^{[-\mathbf{v}_k]_+}]\right)
\end{aligned}
\]
in $\C \otimes K(\mathscr{C})_{\Z}$.
If $h_k > 0$, then $\mathbf{b}(s)_k > 0$ for any
$1 \leq s \leq \ell$ since
$\mathbf{b}(s) - \mathbf{h} \in \Z\mathbf{v}_k$
and $(\mathbf{v}_k)_k = 0$.
It contradicts
\[
\min_{1 \leq s \leq \ell}\mathbf{b}(s)_k = 0.
\]
Thus $[L_k] \nmid [L^{\mathbf{h}}]$ in
$\C \otimes K(\mathscr{C})_{\Z}$ by
Lemma~\ref{Lem : real prime modules give prime elements over C}.
Therefore
\[
[L_k] \mid
\left([L^{[\mathbf{v}_k]_+}]
- \alpha_j[L^{[-\mathbf{v}_k]_+}]\right)
\]
in $\C \otimes K(\mathscr{C})_{\Z}$ for some
$1 \leq j \leq \mathfrak{d}(L_k, S)$.
Now the claim follows by
Lemma~\ref{Lem : Lambda admissibility and binomial divisibility}.
\end{proof}

\begin{Cor}
\label{Cor : complete Lambda admissibility of maximal real sets}
Let $\mathcal{M} = \{L_i \mid i \in J\}$ be a maximal real set.
Then $(\mathcal{M}, \widetilde{B}_{\mathcal{M}})$ is a completely
$\Lambda$-admissible monoidal seed.
\end{Cor}

\begin{proof}
By Theorem~\ref{Thm : Lambda admissibility of maximal real sets},
it is enough to prove that
\[
\mu_k\widetilde{B}_{\mathcal{M}}
= \widetilde{B}_{\mu_k\mathcal{M}}
\qquad\text{for any }k \in J_{\mathrm{uf}}.
\]
Here
$\mu_k\mathcal{M}
= \{L_i \mid i \in J \setminus \{k\}\} \cup \{L'_k\}$
is the maximal real set with
\[
[L'_k] =
\frac{[L^{[\mathbf{v}_k]_+}] + [L^{[-\mathbf{v}_k]_+}]}{[L_k]}.
\]
Since
\[
\frac{1}{2}
\left(\Lambda_{\mu_k\mathcal{M}}
- \Lambda^{\infty}_{\mu_k\mathcal{M}}\right)
\in \mathrm{GL}_N(\Z),
\]
$\widetilde{B}_{\mu_k\mathcal{M}}$ is characterized by
\[
\widetilde{B}_{\mu_k\mathcal{M}}^T\Lambda_{\mu_k\mathcal{M}}
= (2I_{|J_{\mathrm{uf}}|}\ 0),
\qquad
\Lambda^{\infty}_{\mu_k\mathcal{M}}
\widetilde{B}_{\mu_k\mathcal{M}} = 0.
\]
Let
$\widetilde{B}_{\mathcal{M}}
= (b_{ij})_{i \in J,\, j \in J_{\mathrm{uf}}}$.
As in the proof of \cite[Proposition 6.4]{KKOP20},
we can prove that
\[
\begin{aligned}
\Lambda(L_i, L'_k)
&= -\Lambda(L_i, L_k)
+ \sum_{\ell \in J}\Lambda(L_i, L_\ell)[-b_{\ell k}]_+, \\
\Lambda(L'_k, L_i)
&= -\Lambda(L_k, L_i)
+ \sum_{\ell \in J}\Lambda(L_\ell, L_i)[b_{\ell k}]_+
\end{aligned}
\]
for all $i \neq k$.
Thus we have
$\Lambda_{\mu_k\mathcal{M}} = \mu_k\Lambda_{\mathcal{M}}$.
Therefore
\[
(\mu_k\widetilde{B}_{\mathcal{M}})^T
\Lambda_{\mu_k\mathcal{M}} = (2I_{|J_{\mathrm{uf}}|}\ 0)
\]
holds by the fact in Subsection~\ref{ssec : cluster algebras}.

By \cite[(3.2)]{BFZ05},
\[
\mu_k\widetilde{B}_{\mathcal{M}}
= E_{k,+}^{\widetilde{B}_{\mathcal{M}}}
\widetilde{B}_{\mathcal{M}}
F_{k,+}^{\widetilde{B}_{\mathcal{M}}},
\]
where
\begin{itemize}
    \item
    $E_{k,+}^{\widetilde{B}_{\mathcal{M}}} = (e_{ij})_{i,j \in J}$,
    \[
    e_{ij} =
    \begin{cases}
        \delta_{ij}, & \text{if }j \neq k, \\
        -1, & \text{if }(i, j) = (k, k), \\
        [-b_{ik}]_+, & \text{if }i \neq k,\ j = k,
    \end{cases}
    \]
    \item
    $F_{k,+}^{\widetilde{B}_{\mathcal{M}}}
    = (f_{ij})_{i,j \in J_{\mathrm{uf}}}$,
    \[
    f_{ij} =
    \begin{cases}
        \delta_{ij}, & \text{if }i \neq k, \\
        -1, & \text{if }(i, j) = (k, k), \\
        [b_{kj}]_+, & \text{if }i = k,\ j \neq k.
    \end{cases}
    \]
\end{itemize}
Note that
$\left(E_{k,+}^{\widetilde{B}_{\mathcal{M}}}\right)^2 = I_N$.
By Lemma~\ref{Lem : Lambda infinity tensor products and subquotients},
we have
\[
\begin{aligned}
\Lambda^{\infty}(M, L'_k)
&= -\Lambda^{\infty}(M, L_k)
+ \sum_{\ell \in J}\Lambda^{\infty}(M, L_\ell)[-b_{\ell k}]_+, \\
\Lambda^{\infty}(L'_k, M)
&= -\Lambda^{\infty}(L_k, M)
+ \sum_{\ell \in J}\Lambda^{\infty}(L_\ell, M)[b_{\ell k}]_+
\end{aligned}
\]
for all $M \in \Irr\mathscr{C}$.
Thus we obtain that
\[
\Lambda^{\infty}_{\mu_k\mathcal{M}}
= \left(E_{k,+}^{\widetilde{B}_{\mathcal{M}}}\right)^T
\Lambda^{\infty}_{\mathcal{M}}
E_{k,+}^{\widetilde{B}_{\mathcal{M}}}.
\]
Then it holds that
\[
\begin{aligned}
\Lambda^{\infty}_{\mu_k\mathcal{M}}
(\mu_k\widetilde{B}_{\mathcal{M}})
&= \left(E_{k,+}^{\widetilde{B}_{\mathcal{M}}}\right)^T
\Lambda^{\infty}_{\mathcal{M}}
\left(E_{k,+}^{\widetilde{B}_{\mathcal{M}}}\right)^2
\widetilde{B}_{\mathcal{M}}
F_{k,+}^{\widetilde{B}_{\mathcal{M}}} \\
&= \left(E_{k,+}^{\widetilde{B}_{\mathcal{M}}}\right)^T
\Lambda^{\infty}_{\mathcal{M}}\widetilde{B}_{\mathcal{M}}
F_{k,+}^{\widetilde{B}_{\mathcal{M}}} \\
&= 0.
\end{aligned}
\]
\end{proof}

